%-------------------------
% Resume in Latex
% Author : Jake Gutierrez
% Based off of: https://github.com/sb2nov/resume
% License : MIT
%------------------------

\documentclass[letterpaper,11pt]{article}

\usepackage[T1]{fontenc}
\usepackage[utf8]{inputenc}
\usepackage{latexsym}
\usepackage[empty]{fullpage}
\usepackage{titlesec}
\usepackage{marvosym}
\usepackage[usenames,dvipsnames]{color}
\usepackage{verbatim}
\usepackage{enumitem}
\usepackage[hidelinks]{hyperref}
\usepackage{fancyhdr}
\usepackage[english]{babel}
\usepackage{tabularx}
\input{glyphtounicode}


%----------FONT OPTIONS----------
% sans-serif
% \usepackage[sfdefault]{FiraSans}
% \usepackage[sfdefault]{roboto}
% \usepackage[sfdefault]{noto-sans}
% \usepackage[default]{sourcesanspro}

% serif
% \usepackage{CormorantGaramond}
% \usepackage{charter}


\pagestyle{fancy}
\fancyhf{} % clear all header and footer fields
\fancyfoot{}
\renewcommand{\headrulewidth}{0pt}
\renewcommand{\footrulewidth}{0pt}

% Adjust margins
\addtolength{\oddsidemargin}{-0.5in}
\addtolength{\evensidemargin}{-0.5in}
\addtolength{\textwidth}{1in}
\addtolength{\topmargin}{-.5in}
\addtolength{\textheight}{1.0in}

\urlstyle{same}

\raggedbottom
\raggedright
\setlength{\tabcolsep}{0in}

% Sections formatting
\titleformat{\section}{
  \vspace{-4pt}\scshape\raggedright\large
}{}{0em}{}[\color{black}\titlerule \vspace{-5pt}]

% Ensure that generate pdf is machine readable/ATS parsable
\pdfgentounicode=1

%-------------------------
% Custom commands
\newcommand{\resumeItem}[1]{
  \item\small{
    {#1 \vspace{-2pt}}
  }
}

\newcommand{\resumeSubheading}[4]{
  \vspace{-2pt}\item
    \begin{tabular*}{0.97\textwidth}[t]{l@{\extracolsep{\fill}}r}
      \textbf{#1} & #2 \\
      \textit{\small#3} & \textit{\small #4} \\
    \end{tabular*}\vspace{-7pt}
}

\newcommand{\resumeSubSubheading}[2]{
    \item
    \begin{tabular*}{0.97\textwidth}{l@{\extracolsep{\fill}}r}
      \textit{\small#1} & \textit{\small #2} \\
    \end{tabular*}\vspace{-7pt}
}

\newcommand{\resumeProjectHeading}[2]{
    \item
    \begin{tabular*}{0.97\textwidth}{l@{\extracolsep{\fill}}r}
      \small#1 & #2 \\
    \end{tabular*}\vspace{-7pt}
}

\newcommand{\resumeSubItem}[1]{\resumeItem{#1}\vspace{-4pt}}

\renewcommand\labelitemii{$\vcenter{\hbox{\tiny$\bullet$}}$}

\newcommand{\resumeSubHeadingListStart}{\begin{itemize}[leftmargin=0.15in, label={}]}
\newcommand{\resumeSubHeadingListEnd}{\end{itemize}}
\newcommand{\resumeItemListStart}{\begin{itemize}}
\newcommand{\resumeItemListEnd}{\end{itemize}\vspace{-5pt}}

%-------------------------------------------
%%%%%%  RESUME STARTS HERE  %%%%%%%%%%%%%%%%%%%%%%%%%%%%


\begin{document}

%----------HEADING----------
% \begin{tabular*}{\textwidth}{l@{\extracolsep{\fill}}r}
%   \textbf{\href{http://sourabhbajaj.com/}{\Large Sourabh Bajaj}} & Email : \href{mailto:sourabh@sourabhbajaj.com}{sourabh@sourabhbajaj.com}\\
%   \href{http://sourabhbajaj.com/}{http://www.sourabhbajaj.com} & Mobile : +1-123-456-7890 \\
% \end{tabular*}

\begin{center}
    \textbf{\Huge \scshape Longye Chen} \\ \vspace{3pt}
    \textit{Candidature bénévolat -- Science Expérience} \\ \vspace{3pt}
    \small \href{tel:+33602536119}{\underline{+33 6 02 53 61 19}} $|$
    \href{mailto:lchen031205@gmail.com}{\underline{lchen031205@gmail.com}} $|$
    \href{https://www.linkedin.com/in/longye-chen-11ba56309/}{\underline{linkedin.com/in/longye-chen}} $|$
    \href{https://github.com/CHENLongyeLucien}{\underline{github.com/CHENLongyeLucien}}
\end{center}


%-----------EDUCATION-----------
\section{Formation}
  \resumeSubHeadingListStart
    \resumeSubheading
      {Sorbonne Université}{Paris, France}
      {Licence Informatique}{Présent}
    \resumeSubheading
      {Sorbonne Université}{Paris, France}
      {Licence parcours Mathématiques et Informatique}{Sept. 2023 -- Juin 2025}
  \resumeSubHeadingListEnd


%-----------EXPERIENCE-----------
\section{Expérience}
  \resumeSubHeadingListStart

    \resumeSubheading
      {Stagiaire Systèmes, Réseaux et Infrastructure}{Juin 2026 -- Sept. 2026}
      {Lessonia}{Bretagne, France}
      \resumeItemListStart
        \resumeItem{\textbf{Arrêt automatique sur coupure électrique} : benchmark de 3 solutions (Eaton IPP, NUT, Eaton Brightlayer), intégration du cluster Proxmox VE et du SAN DataCore à la supervision et règles d'arrêt ordonné selon les dépendances. \textit{Résultat : sécurité critique perdue lors de la migration VMware vers Proxmox rétablie, infrastructure protégée contre la perte de données et alertes en temps réel}}
        \resumeItem{\textbf{Agent IA auto-hébergé} (Ollama, Open WebUI) : webhooks Proxmox VE et GLPI, outil Python en lecture seule vers la base PostgreSQL métier. \textit{Résultat : synthèses d'incidents générées automatiquement chaque jour et données de l'entreprise interrogeables en langage naturel, en toute sécurité}}
        \resumeItem{\textbf{Administration des utilisateurs} de bout en bout (Active Directory, postes, Kerio Connect, droits d'accès). \textit{Résultat : arrivées et départs pris en charge de façon autonome, accès cohérents et sécurisés}}
        \resumeItem{\textbf{Déploiement d'une mise à jour} sur tout le parc et support quotidien via GLPI. \textit{Résultat : parc homogène, interruption de service minimale et incidents utilisateurs résolus}}
\resumeItemListEnd

  \resumeSubHeadingListEnd


%-----------PROJECTS-----------
\section{Projets}
    \resumeSubHeadingListStart
      \resumeProjectHeading
          {\textbf{\href{https://github.com/CHENLongyeLucien/cinema-finder-poc}{\underline{Cinema Finder}}} -- Simulation Datacom $|$ \emph{React, Leaflet, MapLibre, Material UI}}{Sept. 2026}
          \resumeItemListStart
            \resumeItem{\textbf{Immersion dans une équipe de développement} : simulation de la méthode de travail des équipes Datacom (lecture de la documentation, revue de code, soumission pour relecture). \textit{Résultat : prise en main rapide d'une base de code inconnue et des pratiques professionnelles de revue}}
            \resumeItem{\textbf{Audit d'une application web} de cartographie de cinémas (Australie, Nouvelle-Zélande) et planification d'améliorations. \textit{Résultat : axes d'amélioration identifiés et documentés pour la suite du projet}}
            \resumeItem{\textbf{Débogage} avec les outils de développement du navigateur : identification de la cause racine de 4 bugs (coordonnées inversées sur MapLibre, distance de 0 km masquée, durée des notifications ignorée, cache mémoire non borné) et correctifs ciblés. \textit{Résultat : navigation sur la carte fiabilisée, affichage corrigé et fuite mémoire évitée}}
          \resumeItemListEnd
    \resumeSubHeadingListEnd



%
%-----------PROGRAMMING SKILLS-----------
\section{Technical Skills}
 \begin{itemize}[leftmargin=0.15in, label={}]
    \small{\item{
     \textbf{Languages}{: Java, Python, C/C++, SQL (Postgres), JavaScript, HTML/CSS, R} \\
     \textbf{Frameworks}{: React, Node.js, Flask, JUnit, WordPress, Material-UI, FastAPI} \\
     \textbf{Developer Tools}{: Git, Docker, TravisCI, Google Cloud Platform, VS Code, Visual Studio, PyCharm, IntelliJ, Eclipse} \\
     \textbf{Libraries}{: pandas, NumPy, Matplotlib}
    }}
 \end{itemize}


%-------------------------------------------
\end{document}
